\documentclass[12pt,a4paper]{article}

% ─── Packages ────────────────────────────────────────────────
\usepackage{amsmath,amssymb}
\usepackage{booktabs}
\usepackage[margin=1in]{geometry}
\usepackage{setspace}
\usepackage{hyperref}
\usepackage{xcolor}
\usepackage{float}
\usepackage{fontspec}
\usepackage{xeCJK}
\setCJKmainfont{PingFang TC}
\usepackage{enumitem}
\usepackage{longtable}
\usepackage{array}

\hypersetup{colorlinks=true, linkcolor=blue, citecolor=blue, urlcolor=blue}
\onehalfspacing

% Severity colors
\newcommand{\HIGH}{\textcolor{red}{\textbf{HIGH}}}
\newcommand{\MEDIUM}{\textcolor{orange}{\textbf{MEDIUM}}}
\newcommand{\LOW}{\textcolor{teal}{\textbf{LOW}}}

\title{Academic Review Report (v1)\\
\large ``Can Anything Beat VIX? A Systematic Out-of-Sample Evaluation\\of Eleven Signal Families for Equity Volatility Forecasting\\and Volatility Timing''\\[6pt]
\normalsize Yi-Hao Lai \& VolPred Research System (March 2026)}

\author{Reviewer: Claude Opus 4.6 (Academic Review Agent)}
\date{Review Date: March 30, 2026}

\begin{document}
\maketitle
\tableofcontents
\newpage

%=================================================================
\section{Executive Summary}
%=================================================================

This paper conducts a systematic out-of-sample horse race among eleven signal families for equity volatility forecasting and volatility timing, concluding that no signal can beat the VIX. The paper is well-structured, methodologically rigorous, and addresses an important question. The null result is genuinely informative and the ``drawdown insurance'' reframing is a valuable contribution.

However, several issues ranging from methodological gaps to missing robustness checks reduce the paper's persuasiveness and journal-readiness. I identify \textbf{5 HIGH}, \textbf{8 MEDIUM}, and \textbf{7 LOW} severity issues across 10 review dimensions.

\medskip
\noindent\textbf{Overall Assessment}: \textbf{Major Revision} required. The core finding is sound, but the empirical implementation, statistical testing, and exposition need strengthening before submission to a top-tier journal (JBF, JFE, JFQA).

%=================================================================
\section{Dimension 1: Research Question and Contribution}
\label{sec:d1}
%=================================================================

\subsection*{Strengths}
\begin{itemize}
    \item Clear, well-motivated research question with practical relevance
    \item Three distinct contributions: (i) comprehensive horse race, (ii) drawdown insurance framework, (iii) research frontier delineation
    \item The ``informative null'' framing is original and persuasive
\end{itemize}

\subsection*{Issues}

\begin{enumerate}[label=D1.\arabic*]

\item \MEDIUM\ \textbf{Overclaiming ``first comprehensive horse race''}

The paper claims to conduct ``the first comprehensive, pre-registered horse race'' (p.~2). However:
\begin{itemize}
    \item The study is not formally pre-registered (no OSF/AsPredicted registration number is provided)
    \item The term ``pre-specified'' is more accurate than ``pre-registered,'' which has a specific meaning in empirical research
    \item Competitor studies such as Bali and Zhou (2016), Gonzalez-Perez (2015), and others have conducted multi-signal comparisons, though less systematically
\end{itemize}
\textbf{Recommendation}: Replace ``pre-registered'' with ``pre-specified'' throughout, and add a footnote acknowledging prior multi-signal studies while explaining how this study differs.

\item \LOW\ \textbf{Title length}

At 24 words, the title is excessively long for a journal submission. Most top finance journals prefer titles under 15 words.

\textbf{Recommendation}: Consider ``Can Anything Beat VIX? An Out-of-Sample Evaluation of Volatility Signals'' (12 words).

\end{enumerate}


%=================================================================
\section{Dimension 2: Literature Review and Positioning}
\label{sec:d2}
%=================================================================

\subsection*{Strengths}
\begin{itemize}
    \item Good coverage of the core VIX forecasting literature (Christensen \& Prabhala, Jiang \& Tian, Bekaert \& Hoerova)
    \item Appropriate positioning relative to Moreira \& Muir (2017) volatility-managed portfolios
\end{itemize}

\subsection*{Issues}

\begin{enumerate}[label=D2.\arabic*]

\item \HIGH\ \textbf{Missing key literature on VIX forecasting and volatility horse races}

Several important papers are not cited or discussed:
\begin{itemize}
    \item \textbf{Patton and Sheppard (2015)}: ``Good Volatility, Bad Volatility'' -- directly relevant to volatility forecasting decomposition (cited Bollerslev et al.\ 2020 instead, which is about cross-section of returns)
    \item \textbf{Hansen and Lunde (2005)}: ``A Forecast Comparison of Volatility Models'' -- the canonical volatility model horse race paper, methodologically foundational
    \item \textbf{Poon and Granger (2003)}: ``Forecasting Volatility in Financial Markets: A Review'' -- the definitive survey; conspicuously absent
    \item \textbf{Andersen et al.\ (2003)}: ``Modeling and Forecasting Realized Volatility'' -- foundational for realized volatility methodology used throughout
    \item \textbf{Bali and Zhou (2016)}: ``Risk, Uncertainty, and Expected Returns'' -- tests VRP with multiple signals
    \item \textbf{Baele et al.\ (2019)}: ``Volatility-Managed Portfolios'' comment/extension
    \item \textbf{Gonzalez-Perez (2015)}: Model-free implied volatility forecasting -- a close competitor paper
\end{itemize}
\textbf{Recommendation}: Add a dedicated ``Related Literature'' subsection (Section 2.4 or expand 2.5) covering volatility forecasting horse races, model-free implied volatility studies, and the VRP-return prediction debate.

\item \MEDIUM\ \textbf{Incomplete discussion of Moreira \& Muir (2017) debate}

The paper cites Moreira \& Muir (2017) as motivation but does not discuss the extensive debate that followed:
\begin{itemize}
    \item Cederburg et al.\ (2020): Challenges the Moreira-Muir finding with longer samples
    \item Liu et al.\ (2019): Shows volatility timing works only for some factor portfolios
    \item Barroso and Detzel (2021): Examines economic value of volatility timing
\end{itemize}
\textbf{Recommendation}: Add 1--2 paragraphs on the post-Moreira-Muir debate in Section 1 or Section 2.

\end{enumerate}


%=================================================================
\section{Dimension 3: Methodology and Research Design}
\label{sec:d3}
%=================================================================

\subsection*{Strengths}
\begin{itemize}
    \item Unified evaluation pipeline is a genuine methodological contribution
    \item Strict lag enforcement with code-level implementation (\texttt{signal.shift(1)}) is commendable
    \item Transaction cost accounting is standard but well-implemented
    \item Cross-era validation across five non-overlapping periods is sound
\end{itemize}

\subsection*{Issues}

\begin{enumerate}[label=D3.\arabic*]

\item \HIGH\ \textbf{$R^2_{\text{OOS}}$ definition is non-standard}

Equation (5) defines OOS $R^2$ relative to VIX-only forecasts rather than the historical mean. While the paper acknowledges this (``note that this is defined relative to VIX (not the historical mean)''), this is non-standard and creates two problems:
\begin{itemize}
    \item It conflates ``beating VIX'' with ``forecasting ability'' -- a signal could have strong absolute forecasting power but zero incremental power over VIX
    \item It makes results incomparable to the vast literature that uses the Campbell-Thompson (2008) convention ($R^2_{\text{OOS}}$ relative to the prevailing mean)
    \item Negative OOS $R^2$ relative to VIX does not mean the signal is useless -- it means the signal is worse than VIX, which is a different (and less informative) statement
\end{itemize}
\textbf{Recommendation}: Report both the standard OOS $R^2$ (relative to historical mean) and the VIX-relative OOS $R^2$. Add the Campbell-Thompson reference. This would also show that many signals \emph{do} forecast volatility, just not better than VIX -- strengthening the sufficiency argument.

\item \HIGH\ \textbf{Uneven sample periods across signal families}

The eleven signals have wildly different sample sizes (Table 1): from 811 days (Families 5--6) to 8,325 days (Family 11). This creates serious comparability issues:
\begin{itemize}
    \item Families 5--6 (811 days, 2023--2026) are tested on a period that includes the 2023--2024 bull market and the 2025 tariff volatility, but not GFC, dot-com, or COVID
    \item The DM test power varies dramatically: with 811 observations, a DM test can only detect very large differences; with 8,325, it can detect subtle ones
    \item Era stability analysis (Table 5) cannot include Families 5--6 in Eras 1--4
\end{itemize}
\textbf{Recommendation}: Either (a) restrict the horse race to signals available for the full 1993--2026 period, or (b) add a ``common sample'' analysis where all available signals are compared over 2015--2026 (where most signals overlap), or (c) explicitly acknowledge the power heterogeneity and adjust the Harvey threshold by sample size.

\item \MEDIUM\ \textbf{Rolling window length not justified}

The paper uses a 252-day rolling window for parameter estimation without justification. This choice can affect results materially:
\begin{itemize}
    \item Too short: noisy estimates, frequent regime-change artifacts
    \item Too long: slow adaptation to structural breaks
    \item Sensitivity to window length is not tested
\end{itemize}
\textbf{Recommendation}: Add a robustness check with 126-day and 504-day windows.

\item \MEDIUM\ \textbf{Forecast target mismatch with VIX definition}

The paper forecasts 22-day realized volatility (Equation 3) but VIX targets 30-\emph{calendar}-day (approximately 21--22 trading days) implied volatility. While the approximation is standard, the mismatch should be discussed, especially since small differences compound over thousands of observations.

\textbf{Recommendation}: Add a sentence acknowledging the VIX targets 30 calendar days ($\approx$ 21 trading days) and explain why 22 trading days is used.

\end{enumerate}


%=================================================================
\section{Dimension 4: Statistical Testing and Inference}
\label{sec:d4}
%=================================================================

\subsection*{Strengths}
\begin{itemize}
    \item Harvey (2016) $|t|>3.0$ threshold is appropriate for multiple testing
    \item Diebold-Mariano test with Newey-West standard errors is correct for multi-step forecasts
\end{itemize}

\subsection*{Issues}

\begin{enumerate}[label=D4.\arabic*]

\item \HIGH\ \textbf{No formal multiple-testing correction procedure}

The paper adopts $|t|>3.0$ as a rule of thumb from Harvey et al.\ (2016) but does not implement a formal multiple-testing procedure:
\begin{itemize}
    \item Harvey et al.\ (2016) recommend a \emph{Bayesian} adjustment or a haircut formula ($t_{\text{adj}} = t \times \sqrt{1 - \pi_0}$), not simply raising the threshold to 3.0
    \item With 11 tests, a Bonferroni correction would set the threshold at $|t| > 3.23$ (using $\alpha = 0.05/11$); the Holm-Bonferroni or BH FDR procedures would be more powerful
    \item The $|t|>3.0$ threshold is used for both volatility forecasting (Table 2) and strategy comparison (Table 3) -- but these are different hypothesis families with different multiplicity concerns
    \item The paper does not report p-values, making it impossible for readers to apply their own corrections
\end{itemize}
\textbf{Recommendation}: (a) Implement either Holm-Bonferroni or Benjamini-Hochberg FDR across the 11 DM tests and report adjusted p-values; (b) report raw p-values for all tests; (c) clarify that $|t|>3.0$ is an approximate threshold from Harvey et al.\ but also show formal corrections. The conclusion would almost certainly be the same (all signals fail), but the methodology would be ironclad.

\item \MEDIUM\ \textbf{No bootstrap or permutation-based inference}

The paper relies entirely on asymptotic inference (DM test, t-tests). Given the overlapping nature of 22-day realized volatility observations, the effective sample size is much smaller than $T$:
\begin{itemize}
    \item 22-day RV creates $\sim 22$-day overlap, reducing effective degrees of freedom by a factor of $\sim 22$
    \item Newey-West corrections address serial correlation but may be insufficient for the degree of overlap
    \item Block bootstrap or circular bootstrap would provide more reliable inference
\end{itemize}
\textbf{Recommendation}: Add a bootstrap robustness check (block bootstrap with block length 22--44 days) for the main DM tests.

\item \LOW\ \textbf{Significance stars inconsistent with Harvey threshold}

Table 2 uses standard significance stars ($^{***}$, $^{**}$, $^{*}$) for in-sample t-statistics while simultaneously advocating $|t|>3.0$. This sends a mixed message: the VIX term structure has $t=17.6^{***}$ in-sample but DM $|t|=0.87$ out-of-sample. The stars on in-sample statistics may mislead readers into thinking these are ``significant'' in the Harvey sense.

\textbf{Recommendation}: Either remove significance stars from in-sample statistics or add a clear note: ``In-sample stars use conventional thresholds; all out-of-sample tests require $|t|>3.0$.''

\end{enumerate}


%=================================================================
\section{Dimension 5: Empirical Results and Tables}
\label{sec:d5}
%=================================================================

\subsection*{Strengths}
\begin{itemize}
    \item Eight well-formatted tables covering the main results, era stability, and robustness
    \item The cross-era stability analysis (Table 5) is particularly well-done
    \item The insurance framework table (Table 8) is original and useful
\end{itemize}

\subsection*{Issues}

\begin{enumerate}[label=D5.\arabic*]

\item \MEDIUM\ \textbf{Table 2: Missing or incomplete results for Families 5, 6, 7}

Three of eleven signal families have incomplete results in Table 2:
\begin{itemize}
    \item Families 5 (Multi-asset optimization) and 6 (ERC): marked as ``---'' because they ``are portfolio strategies without direct volatility forecasting components.'' But the paper's title promises evaluation of eleven families -- presenting only 8--9 complete rows undermines the horse race claim
    \item Family 7 (Bitcoin volatility): partial correlation reported (0.178) but all other columns are ``---''
\end{itemize}
\textbf{Recommendation}: Either (a) complete the analysis for all 11 families (even portfolio strategies can be evaluated on portfolio volatility forecasting), or (b) clearly separate the paper into two parts: a ``volatility forecasting'' horse race (Families 1--4, 7--11) and a ``portfolio strategy'' comparison (Families 5--6), and adjust the title/abstract accordingly.

\item \MEDIUM\ \textbf{Table 3: Inconsistent strategy numbering}

Table 3 lists strategies 1, 2, 3, 4, 4, 6, 8, 9, 10, 11. Family 4 appears twice (VRP timing and VRP percentile), Family 5 is missing (covered in Table 4), and Family 7 is missing entirely. This makes cross-referencing Tables 2 and 3 difficult.

\textbf{Recommendation}: Use consistent numbering; label the two VRP strategies as ``4a'' and ``4b''; add a row for Family 7 (Bitcoin) even if it underperforms.

\item \LOW\ \textbf{Conflicting Sharpe ratios for the 50/50 baseline}

The paper reports two different Sharpe ratios for the 50/50 SPY/GLD allocation:
\begin{itemize}
    \item Table 3: Sharpe = 0.947 (long sample)
    \item Table 4: Sharpe = 1.849 (2023-01-04 to 2026-03-27, 811 days)
\end{itemize}
While these reflect different sample periods, the nearly 2$\times$ difference is jarring and could confuse readers. The Table 4 period (2023--2026) is a historically unusual bull market; reporting Sharpe = 1.849 as the benchmark for multi-asset optimization may be misleading.

\textbf{Recommendation}: Add a footnote to Table 4 explicitly noting the different sample period and cautioning against cross-table comparison.

\item \LOW\ \textbf{No figures}

The paper has zero figures. For a 31-page paper with 8 tables, the absence of any visual illustration is notable:
\begin{itemize}
    \item A time-series plot of VIX vs.\ 22-day RV across the full sample would illustrate the VIX--RV relationship
    \item A bar chart of OOS $R^2$ or DM statistics across families would make the null result visually striking
    \item The era-by-era results would benefit from a panel plot
    \item The drawdown insurance trade-off could be shown as a $\gamma$ vs.\ welfare curve
\end{itemize}
\textbf{Recommendation}: Add 3--4 figures: (i) VIX vs.\ RV time series, (ii) DM statistics bar chart, (iii) era stability panel, (iv) welfare as a function of $\gamma$.

\end{enumerate}


%=================================================================
\section{Dimension 6: Theoretical Framework}
\label{sec:d6}
%=================================================================

\subsection*{Strengths}
\begin{itemize}
    \item The information aggregation argument (Section 8.1) is intuitive and well-presented
    \item The VIX formula (Equation 1) and VRP decomposition (Equation 2) are correctly stated
    \item Proposition 1 (break-even $\gamma^*$) provides a clean analytical framework
\end{itemize}

\subsection*{Issues}

\begin{enumerate}[label=D6.\arabic*]

\item \MEDIUM\ \textbf{Proposition 1 derivation is incomplete}

The break-even risk aversion formula (Equation 12) appears to use a second-order approximation of CRRA utility:
$$\gamma^* = \frac{\Delta \mu}{\frac{1}{2}(\sigma_{\text{BH}}^2 - \sigma_{\text{VT}}^2)}$$
But this approximation is only valid for small $\gamma$ and normally distributed returns. For $\gamma = 4.5$ and the fat-tailed returns typical of equity markets, the approximation error can be substantial. The paper does not:
\begin{itemize}
    \item State the approximation assumptions
    \item Compare the approximate $\gamma^*$ to the exact numerical solution
    \item Discuss the impact of skewness and kurtosis on the break-even calculation
\end{itemize}
\textbf{Recommendation}: Either (a) derive $\gamma^*$ from the exact CRRA certainty-equivalent formula using numerical integration, or (b) add a footnote acknowledging the second-order approximation and reporting the exact $\gamma^*$ from simulation.

\item \LOW\ \textbf{Sufficiency claim requires stronger conditions}

Equation (9) claims VIX is a ``sufficient statistic'' for $\sigma^2_{t+1:t+22}$ conditional on $\Omega_t$. This is a strong claim that requires:
\begin{itemize}
    \item Semi-strong efficiency of the options market (stated but not tested)
    \item All signals $s_t$ being elements of $\Omega_t$ (claimed but some signals like Google Trends may not be actively monitored by option traders)
    \item VIX being a function only of option prices (true by construction, but VIX methodology changes in 2003 are not discussed)
\end{itemize}
\textbf{Recommendation}: Soften the language from ``sufficient statistic'' (which has a precise statistical meaning) to ``approximately sufficient conditional predictor'' and add a caveat about signals outside $\Omega_t$.

\end{enumerate}


%=================================================================
\section{Dimension 7: Robustness and Sensitivity}
\label{sec:d7}
%=================================================================

\subsection*{Strengths}
\begin{itemize}
    \item Five-era cross-validation provides regime robustness
    \item Transaction cost sensitivity is implicitly tested (5 bps per leg is reasonable)
\end{itemize}

\subsection*{Issues}

\begin{enumerate}[label=D7.\arabic*]

\item \HIGH\ \textbf{Missing key robustness checks}

Several standard robustness checks are absent:
\begin{itemize}
    \item \textbf{Forecast horizon}: Only 22-day RV is tested. VIX forecasts 30-calendar-day volatility, but practitioners also care about 5-day (weekly) and 63-day (quarterly) horizons. VIX sufficiency may break at shorter horizons where intraday information matters more.
    \item \textbf{Alternative loss functions}: Only squared errors (for DM test) are used. QLIKE loss is standard in volatility forecasting and is more robust to outliers (Patton, 2011). The 41.8\% QLIKE improvement for HAR-RV mentioned in the text should be presented systematically.
    \item \textbf{Parameter stability}: No rolling window sensitivity (126, 252, 504 days).
    \item \textbf{Conditional analysis}: VIX sufficiency may hold on average but break in specific regimes (e.g., during VIX $> 30$, signals might add information about the \emph{persistence} of high volatility).
    \item \textbf{International markets}: Only S\&P 500 is tested. VIX sufficiency for non-U.S.\ markets (e.g., using VSTOXX for European equities, VHSI for Hong Kong, VIXTWN for Taiwan) would strengthen generalizability.
\end{itemize}
\textbf{Recommendation}: Add a dedicated robustness section (even if brief) covering at least forecast horizon sensitivity and QLIKE loss. International evidence would strengthen the paper substantially.

\end{enumerate}


%=================================================================
\section{Dimension 8: Writing Quality and Exposition}
\label{sec:d8}
%=================================================================

\subsection*{Strengths}
\begin{itemize}
    \item Clear, well-organized prose with logical flow
    \item Good use of numbered equations and cross-references
    \item The five-point methodological safeguard list in the introduction is effective
\end{itemize}

\subsection*{Issues}

\begin{enumerate}[label=D8.\arabic*]

\item \LOW\ \textbf{Self-referential footnotes reduce credibility}

Footnote 1 (p.~2) references the authors' own ``research program'' and ``Codex code review'' of specific experiment IDs (K618, K621, K679, K698). While honest and transparent, this is unusual for a journal article and may signal to reviewers that the paper is a by-product of a software project rather than an independent academic contribution.

\textbf{Recommendation}: Rephrase the footnote to describe the finding generically: ``In preliminary analyses, we identified four instances where look-ahead bias inflated Sharpe ratios by up to 373\%, underscoring the importance of strict lag enforcement.''

\item \LOW\ \textbf{Section 8.4 ``Simplicity Premium'' is underspecified}

The meta-analysis of 14 live-tracked strategies is mentioned but the strategies are not listed, the tracking period is not specified, and the correlation ($\rho = 0.077$, $p = 0.794$) is based on only 14 observations, which is too small for meaningful inference.

\textbf{Recommendation}: Either expand Section 8.4 with a table of the 14 strategies and their characteristics, or move this result to a footnote and tone down the claim.

\item \LOW\ \textbf{Inconsistent notation}

Minor notation inconsistencies:
\begin{itemize}
    \item $\text{VIX}_t$ vs.\ $\text{VIX}_{t-1}$: Equation (8) uses $\text{VIX}_{t-1}$ but Equations (4)--(5) use $\text{VIX}_t$. The lag convention should be consistent.
    \item $RV_{t+1:t+22}$ vs.\ $RV_t^{(22)}$: Both are used; clarify whether $RV_t^{(22)}$ is backward-looking (past 22 days ending at $t$) while $RV_{t+1:t+22}$ is forward-looking.
\end{itemize}
\textbf{Recommendation}: Add a notation table or a paragraph at the start of Section 3 clarifying all RV and VIX subscript conventions.

\end{enumerate}


%=================================================================
\section{Dimension 9: Originality and Novelty}
\label{sec:d9}
%=================================================================

\subsection*{Strengths}
\begin{itemize}
    \item The systematic multi-signal horse race with unified methodology is valuable
    \item The ``informative null'' framing is novel and publishable
    \item The drawdown insurance framework connects VT to utility theory
\end{itemize}

\subsection*{Issues}

\begin{enumerate}[label=D9.\arabic*]

\item \MEDIUM\ \textbf{Limited methodological novelty}

The paper's methodology (regression, DM test, rolling window) is entirely standard. The contribution is empirical (testing more signals systematically) rather than methodological. This positions the paper well for JBF but may face resistance at JFE or RFS, which typically require methodological innovation.

The paper would benefit from at least one novel methodological element:
\begin{itemize}
    \item A model confidence set (MCS; Hansen et al., 2011) to formally rank all signals simultaneously
    \item A Bayesian model averaging approach to optimally combine signals
    \item A conditional sufficiency test (is VIX sufficient in all regimes, or only some?)
\end{itemize}
\textbf{Recommendation}: Add an MCS analysis. The MCS would likely contain only VIX, providing a clean formal confirmation of the paper's thesis. The MCS is computationally straightforward and would add a genuine methodological contribution.

\end{enumerate}


%=================================================================
\section{Dimension 10: Journal Fit and Submission Readiness}
\label{sec:d10}
%=================================================================

\subsection*{Assessment}

\begin{itemize}
    \item \textbf{Target: Journal of Banking \& Finance (JBF)}: Good fit. JBF publishes comprehensive empirical studies with practical relevance. The multiple-signal horse race format and the null result with practical implications align well.
    \item \textbf{Alternative: Journal of Financial and Quantitative Analysis (JFQA)}: Possible but harder sell. Would require stronger methodological contribution (MCS, conditional testing).
    \item \textbf{Journal of Financial Econometrics (JoFE)}: Good fit if the paper adds HAR-RV comparisons and QLIKE loss.
    \item \textbf{Journal of Financial Economics (JFE) / Review of Financial Studies (RFS)}: Unlikely without a theoretical model or striking positive finding.
\end{itemize}

\subsection*{Issues}

\begin{enumerate}[label=D10.\arabic*]

\item \MEDIUM\ \textbf{Paper length and structure need tightening}

At 31 pages with 8 tables, the paper is somewhat long for its content. Several sections can be condensed:
\begin{itemize}
    \item Section 2 (Why VIX is the Benchmark) overlaps with the Introduction; merge into a shorter Section 2 or fold into Introduction
    \item Tables 5 and 6 could be combined into a single era-stability table
    \item Section 6 (Strategy Design) is formulaic; condense to 1 page
\end{itemize}
\textbf{Recommendation}: Target 25--28 pages with tighter exposition and combined tables, plus a supplementary appendix for detailed signal construction.

\end{enumerate}


%=================================================================
\section{Summary of Issues by Severity}
\label{sec:summary}
%=================================================================

\begin{longtable}{@{}p{1cm}p{2cm}p{10cm}@{}}
\toprule
\textbf{ID} & \textbf{Severity} & \textbf{Description} \\
\midrule
\endfirsthead
\toprule
\textbf{ID} & \textbf{Severity} & \textbf{Description} \\
\midrule
\endhead
\bottomrule
\endlastfoot
D2.1 & \HIGH & Missing key literature (Hansen \& Lunde 2005, Poon \& Granger 2003, Andersen et al.\ 2003, Patton \& Sheppard 2015) \\
D3.1 & \HIGH & Non-standard OOS $R^2$ definition (relative to VIX, not historical mean) \\
D3.2 & \HIGH & Uneven sample periods (811 to 8,325 days) with no power analysis \\
D4.1 & \HIGH & No formal multiple-testing correction (only $|t|>3.0$ rule of thumb) \\
D7.1 & \HIGH & Missing key robustness checks (forecast horizon, QLIKE, international, conditional) \\
\midrule
D1.1 & \MEDIUM & Overclaiming ``first/pre-registered'' horse race \\
D2.2 & \MEDIUM & Incomplete discussion of post-Moreira-Muir debate \\
D3.3 & \MEDIUM & Rolling window length (252 days) not justified or sensitivity-tested \\
D3.4 & \MEDIUM & Forecast target (22-day) vs.\ VIX definition (30-calendar-day) mismatch \\
D5.1 & \MEDIUM & Missing/incomplete results for Families 5, 6, 7 in Table 2 \\
D5.2 & \MEDIUM & Inconsistent strategy numbering in Table 3 \\
D6.1 & \MEDIUM & Proposition 1 uses second-order approximation without justification \\
D9.1 & \MEDIUM & Limited methodological novelty (no MCS, no conditional testing) \\
D10.1 & \MEDIUM & Paper length can be tightened (31 pages $\rightarrow$ 25--28) \\
\midrule
D1.2 & \LOW & Title too long (24 words) \\
D4.2 & \LOW & No bootstrap inference for DM tests \\
D4.3 & \LOW & Significance stars inconsistent with Harvey threshold \\
D5.3 & \LOW & Conflicting Sharpe ratios for 50/50 baseline across tables \\
D5.4 & \LOW & No figures in the entire paper \\
D6.2 & \LOW & ``Sufficient statistic'' claim needs softer language \\
D8.1 & \LOW & Self-referential footnotes (internal experiment IDs) \\
D8.2 & \LOW & Section 8.4 Simplicity Premium underspecified ($n=14$) \\
D8.3 & \LOW & Minor notation inconsistencies ($\text{VIX}_t$ vs.\ $\text{VIX}_{t-1}$) \\
\end{longtable}

\noindent\textbf{Total}: 5 HIGH, 9 MEDIUM, 9 LOW = 23 issues.


%=================================================================
\section{Priority Action Items}
\label{sec:actions}
%=================================================================

\noindent\textbf{Must fix before submission (HIGH):}
\begin{enumerate}
    \item Add missing foundational references (Poon \& Granger 2003, Hansen \& Lunde 2005, Andersen et al.\ 2003)
    \item Report standard OOS $R^2$ (relative to historical mean) alongside VIX-relative $R^2$
    \item Address sample size heterogeneity with a common-sample analysis or power discussion
    \item Implement formal multiple-testing correction (BH FDR or Holm-Bonferroni)
    \item Add robustness checks: QLIKE loss, 5-day and 63-day forecast horizons, at least one non-U.S.\ market
\end{enumerate}

\noindent\textbf{Strongly recommended (MEDIUM):}
\begin{enumerate}[resume]
    \item Replace ``pre-registered'' with ``pre-specified''
    \item Add post-Moreira-Muir debate discussion
    \item Justify rolling window length and add sensitivity
    \item Complete Table 2 for all 11 families
    \item Fix Table 3 numbering
    \item Provide exact CRRA calculation for Proposition 1
    \item Add MCS analysis
    \item Tighten paper length
\end{enumerate}

\noindent\textbf{Nice to have (LOW):}
\begin{enumerate}[resume]
    \item Shorten title
    \item Add bootstrap DM tests
    \item Remove/generalize self-referential footnotes
    \item Add 3--4 figures
    \item Add notation table
\end{enumerate}


%=================================================================
\section{Concluding Remarks}
\label{sec:remarks}
%=================================================================

This paper addresses an important and practically relevant question -- whether any signal can improve upon VIX for volatility forecasting and timing -- and delivers a clear, well-supported answer: no, at daily frequency with publicly available data. The ``informative null'' framing is the paper's strongest conceptual contribution, and the drawdown insurance reframing adds practical value.

The main weaknesses are: (i) gaps in the literature review that omit foundational volatility forecasting papers, (ii) methodological choices that are reasonable but not fully justified (non-standard $R^2$, no formal multiple testing, uneven samples), and (iii) missing robustness checks that are now standard in the volatility forecasting literature (QLIKE loss, multiple horizons, international markets).

With the HIGH-severity revisions addressed, this paper would be a strong candidate for the Journal of Banking \& Finance. The null result, properly contextualized, is a genuine contribution to the literature.

\end{document}
